\subsection{Stability of the cube condition under perturbations}\label{ssec:rSstab}

Subsection~\ref{ssec:fffree} proves the averaged cube condition of
Definition~\ref{def:avgcube} for the free massive fields, with a constant of
order $(m\ell)^2$ that depends on the lattice spacing only through the mass in
units of the cube. This subsection asks which perturbations of the Gaussian
preserve it. The answer has a negative and a positive half. A lower bound on
the Hessian of the Hamiltonian does not preserve it: a convex perturbation of
the free field can drive the constant to zero, because the Dirichlet form of
the cube dynamics, unlike the carr\'e du champ of the Langevin dynamics,
depends on the Hessian and is not monotone in it. What preserves it is a
two-sided control of the Hessian, or more precisely a pointwise stable
decomposition of the Hessian along the cubes; under that hypothesis the
condition holds for every strictly log-concave measure that satisfies it, with
a constant that is uniform in the volume and in the cube side, and the proof is
Lions' lemma lifted to the Helffer--Sj\"ostrand representation. Bounded
perturbations are not controlled by their oscillation: the Ising model on the
torus is a perturbation of a product measure with bounded oscillation on every
cube and violates the condition with a constant that vanishes with the volume.
For the lattice gauge field in a fixed gauge the hypotheses of the positive
theorems fail at every coupling, and the subsection records exactly where: the
Hessian of the Wilson action has the flat transverse fields in its kernel at
every configuration, and is not convex on a set of positive measure. The
quadratic approximation of the Wilson action is massless, its massive
regularisation has a cube constant that vanishes with the mass at every
coupling, and the missing estimate is a mass that the generator part of the
Witten Laplacian, not the potential, would have to supply.

\emph{Setting.} We keep the notation of Sec.~\ref{ssec:fffree}: $d\ge1$,
$n\ge2$, $L\ge n+1$, $\mathbb T=(\Z/L\Z)^d$, the coordinate space $\R^N$
(sites or links), the sets $V_x$, $x\in\mathbb T$ (for links, $V_x$ stands for
the set $E_x$ of links with base point in $x+\{0,\dots,n-1\}^d$), each
coordinate lying in exactly $n^d$ of them, the profile $\tau$ with the
constants $S_2$, $S_{11}$, $D_2$, the parabolic constants
$\alpha_n,\beta_n,\rho_n$ of Lemma~\ref{lem:ffparabola}, and the functions
$\chi_x(y)=\prod_i\tau(y_i-x_i)$, which we regard as diagonal operators on
$\R^N$ (for links, $(\chi_xA)(y,i)=\chi_x(y)A(y,i)$). A \emph{partition of
unity} is a family $(\chi_x)_{x\in X}$ of diagonal operators with nonnegative
entries, $\chi_x$ supported in $V_x$ and $\sum_x\chi_x=I$; the family above is
one, and for it $\sum_x\chi_x^2=S_2^dI$ by~\eqref{eq:ffchisums}. For a
positive definite $Q_0$ we write
\begin{equation}\label{eq:rSC0}
C_0(Q_0)=\inf\Bigl\{C:\sum_{x}\ip{\chi_xv}{Q_0\chi_xv}\le C\ip v{Q_0v}
\text{ for all }v\in\R^N\Bigr\},
\end{equation}
the stable decomposition constant of $Q_0$ for the parabolic partition of
unity, and $C_0(Q_0;T)$ for the same infimum over $v$ in a subspace $T$.
Lemma~\ref{lem:ffpartition} gives $C_0(Q_{\rm s})\le S_2^{d-1}(S_2+dD_2/\mu^2)$,
so that $n^dC_0(Q_{\rm s})\le\alpha_n^d(1+10d\rho_ns^{-2})$, $s=n\mu$, and
Lemma~\ref{lem:ffprocadecomp} gives the same bound for $C_0(Q_{\rm P};\ker
\delta_0)$ and $S_2^{d-1}(S_2+\kappa_dD_2/\mu^2)$ for $C_0(Q_{\rm P})$.

A \emph{measure} in this subsection is a probability measure $\pi$ on $\R^N$
with a density $Z^{-1}e^{-H}$, $H\in C^\infty(\R^N)$. For $x\in X$ let
$\Var^\pi_x(f)$ be the conditional variance of $f$ given the coordinates
outside $V_x$, so that $\pi(\Var^\pi_xf)=\lVert f-\mathcal P_xf\rVert^2_{L^2
(\pi)}$, and say that $\pi$ satisfies $\mathrm{BP}(n,\varpi)$ on a class
$\mathcal C\subset L^2(\pi)$ if
\begin{equation}\label{eq:rSBP}
\varpi\Var_\pi(f)\le\frac1{n^d}\sum_{x\in X}\pi\bigl(\Var^\pi_xf\bigr)
\qquad(f\in\mathcal C),
\end{equation}
which for the Wilson measure and $\mathcal C=\mathcal A$ is
Definition~\ref{def:avgcube}. Both sides of~\eqref{eq:rSBP} make sense for
any probability measure on a product of measurable spaces indexed by the
coordinates, and Propositions~\ref{prop:rSglobalHS}, \ref{prop:rSproduct}
and~\ref{prop:rSising} use them in that generality. We write $\varpi(\pi)$
for the largest such $\varpi$ on $\mathcal C=L^2(\pi)$ and $\varpi(Q)$ for
$\varpi(\gamma_Q)$. The conditional law of $\pi$ given the coordinates
outside $V_x$ has the Hamiltonian $\phi_{V_x}\mapsto H(\phi_{V_x},
\phi_{V_x^c})$ on $\R^{V_x}$, whose Hessian is the principal submatrix
$(\Hess H)_{V_xV_x}$.

\subsubsection{A Hessian lower bound does not suffice}

\begin{theorem}[convex perturbations can destroy the condition]
\label{thm:rSnonmono}
Let $\mu_0=\gamma_{Q_{\rm s}}$ be the free scalar field of
Theorem~\ref{thm:ffscalar} and, for $t\ge0$, $\mu_t=Z_t^{-1}e^{-V_t}\mu_0$
with the convex function $V_t(\phi)=\frac t2\lVert\nabla\phi\rVert^2$. Then
$\Hess(-\log\mu_t)=Q_{\rm s}+t(-\Delta)\ge Q_{\rm s}$ for every $t$, while
\[
\varpi(\mu_t)\le\frac{s^2}{s^2+(1+t)\,16d/9}\xrightarrow[t\to\infty]{}0 .
\]
The same holds for the Proca field with $V_t(A)=\frac t2\lVert d_1A\rVert^2$,
on gauge-invariant functions, with the bound $s^2/(s^2+1+t)$. Consequently
there is no $c>0$ such that, for all measures $\mu=e^{-V}\mu_0$ with $V$
convex, $\mathrm{BP}(n,\varpi_0)$ for $\mu_0$ implies $\mathrm{BP}(n,c
\varpi_0)$ for $\mu$; and the same holds with ``$\Hess(-\log\mu)\ge\Hess
(-\log\mu_0)$'' in place of ``$V$ convex''.
\end{theorem}

\begin{proof}
$\mu_t$ is the Gaussian with precision $Q_t=(1+t)(-\Delta)+\mu^2=(1+t)
(-\Delta+\mu_t^2)$, $\mu_t^2=\mu^2/(1+t)$. For $c>0$ the law $\gamma_{cQ}$ is
the image of $\gamma_Q$ under $\phi\mapsto c^{-1/2}\phi$, which commutes with
the conditional expectations $\mathcal P_x$, so $\varpi(cQ)=\varpi(Q)$ and
$\varpi(Q_t)=\varpi(-\Delta+\mu_t^2)$. Proposition~\ref{prop:ffupper}(i) with
$s_t=n\mu_t$ gives $\varpi(Q_t)\le s_t^2/(s_t^2+16d/9)=s^2/(s^2+(1+t)16d/9)$.
The Proca case is Proposition~\ref{prop:ffupper}(ii) in the same way. The
last claim follows because $\mu_0$ itself satisfies $\mathrm{BP}(n,\varpi_0)$
with $\varpi_0=c_d(s)>0$ by Theorem~\ref{thm:ffscalar}, while
$\inf_t\varpi(\mu_t)=0$.
\end{proof}

The mechanism is visible in Lemma~\ref{lem:fflinear}: the Dirichlet form of
the cube dynamics on linear functions is $\ip h{Sh}$ with $S=\sum_xR_x^{\mathsf
T}Q_{V_xV_x}^{-1}R_x$, which \emph{decreases} when $Q$ increases, while the
variance $\ip h{Q^{-1}h}$ also decreases; their ratio has no monotonicity.
For the Langevin dynamics the Dirichlet form is $\pi(\lvert\nabla f\rvert^2)$,
whose integrand does not depend on $\pi$, which is why the Brascamp--Lieb
inequality~\cite{BL76} and the Bakry--\'Emery criterion~\cite{BE85} are
monotone in the Hessian and the cube condition is not. The perturbation $V_t$
is gradient-like: it stiffens the field across the boundaries of the cubes
and so transfers variance to wavelengths longer than the cube, where a cube
resampling cannot reach it. Zeroth-order (mass-like) perturbations behave in
the opposite way, which is the content of the next results.

\subsubsection{Two-sided comparison for Gaussians}

\begin{theorem}[two-sided Gaussian comparison]\label{thm:rSgausscomp}
Let $Q_0$ and $Q$ be positive definite with $Q_0\le Q\le KQ_0$ for some
$K\ge1$. Then $\varpi(Q)\ge\varpi(Q_0)/K$. If $G$ is a subspace compatible
with both $Q_0$ and $Q$, the same holds for the infima in~\eqref{eq:ffvarpi}
over $L^2_G$.
\end{theorem}

\begin{proof}
By Theorem~\ref{thm:ffchaos} and Lemma~\ref{lem:fflinear}, $n^d\varpi(Q)$ is
the infimum over $h\ne0$ of $\sum_x\ip{h_V}{Q_{VV}^{-1}h_V}/\ip h{Q^{-1}h}$,
with $V=V_x$ and $h_V=R_Vh$, and the infimum over $L^2_G$ is the same infimum
over $h\in G^\perp$. Principal submatrices preserve the order, so $Q_{VV}
\le KQ_{0,VV}$ and $Q_{VV}^{-1}\ge K^{-1}Q_{0,VV}^{-1}$; and $Q^{-1}\le
Q_0^{-1}$. Hence each quotient for $Q$ is at least $K^{-1}$ times the
quotient for $Q_0$ at the same $h$.
\end{proof}

\subsubsection{Lions' lemma for the Witten Laplacian}

The Helffer--Sj\"ostrand representation~\cite{HS94,Hel98} expresses the
variance of a strictly log-concave measure through the Witten Laplacian on
one-forms. We use it in the following form. Let $\pi=Z^{-1}e^{-H}d\phi$ with
$H\in C^\infty(\R^N)$ and $\Hess H(\phi)\ge cI$ for all $\phi$ and some
$c>0$. Let $\mathcal L=-\Delta+\ip{\nabla H}\nabla$ be the generator of the
Langevin dynamics, the self-adjoint operator on $L^2(\pi)$ associated with the
Dirichlet form $\pi(\lvert\nabla f\rvert^2)$, and on the space $L^2(\pi;\R^N)$
of vector fields let
\begin{equation}\label{eq:rSwitten}
\mathcal A=\mathcal L\otimes I+\Hess H,
\end{equation}
that is, $\mathcal L$ acting on each component and the multiplication
operator by the matrix $\Hess H(\phi)$; its quadratic form is
\[
a(G,G)=\sum_{i=1}^N\pi\bigl(\lvert\nabla G_i\rvert^2\bigr)+\pi\bigl(\ip G
{\Hess H\,G}\bigr)\ge c\,\pi(\lvert G\rvert^2),
\]
on the form domain $\mathcal D$ of fields whose components have square
integrable weak gradients. Then for every $f\in C^\infty(\R^N)$ with bounded
derivatives of all orders, the class to which the theorem below is first
reduced,
\begin{equation}\label{eq:rShs}
\Var_\pi f=\pi\bigl(\ip{\nabla f}{\mathcal A^{-1}\nabla f}\bigr)=\sup_{G\in
\mathcal D}\bigl[2\pi(\ip{\nabla f}G)-a(G,G)\bigr],
\end{equation}
the supremum being attained at $G=\mathcal A^{-1}\nabla f$. The identity is
the intertwining relation $\nabla\mathcal Lu=\mathcal A\nabla u$ applied to
$u=\mathcal L^{-1}(f-\pi(f))$, which exists because $\mathcal L$ has
spectral gap at least $c$ by the Bakry--\'Emery criterion; the second
equality is the variational characterization of the inverse of a positive
operator. Only two consequences of~\eqref{eq:rShs} are used below: the
identity $\Var_\pi f=2\pi(\ip{\nabla f}{G^*})-a(G^*,G^*)$ at $G^*=\mathcal
A^{-1}\nabla f=\nabla u$, and the inequality $\Var_\pi f\ge2\pi(\ip{\nabla f}
G)-a(G,G)$ for every $G\in\mathcal D$, which follows from $a(\nabla u,G)=\pi
(\ip{\nabla\mathcal Lu}G)=\pi(\ip{\nabla f}G)$ and $a\ge0$ alone, by
expanding $a(\nabla u-G,\nabla u-G)\ge0$. We shall apply~\eqref{eq:rShs} to
$\pi$ and to its conditional laws given the coordinates outside a set $V$,
whose Hamiltonians $\phi_V\mapsto H(\phi_V,\phi_{V^c})$ satisfy the same
hypotheses with the same $c$; for these we write $\mathcal A_V$ and $a_V$, so
that for a field $g$ on $\R^V$,
\[
a_V(g,g)=\sum_{i\in V}\pi_V\bigl(\lvert\nabla_Vg_i\rvert^2\bigr)+\pi_V\bigl(
\ip g{(\Hess H)_{VV}\,g}\bigr),
\]
with $\pi_V$ the conditional law and $\nabla_V$ the gradient in the
coordinates of $V$. Two features of these forms carry the whole argument:
the conditional form differentiates in fewer variables than the global one,
and its Hessian is a principal submatrix of the global Hessian.

\begin{theorem}[the cube condition for log-concave measures]\label{thm:rSwitten}
Let $\pi=Z^{-1}e^{-H}d\phi$ with $H\in C^\infty(\R^N)$ and $\Hess H\ge cI$,
$c>0$. Let $(\chi_x)_{x\in X}$ be a partition of unity for the sets $(V_x)$
and $m_2=\max_i\sum_x\chi_x(i)^2$. Let $T\subset\R^N$ be a subspace such
that $\Hess H(\phi)$ commutes with the orthogonal projection onto $T$ for
every $\phi$ (for instance $T=\R^N$), and suppose that there is $C_H<\infty$
with
\begin{equation}\label{eq:rSstable}
\sum_{x\in X}\ip{\chi_xg}{\Hess H(\phi)\,\chi_xg}\le C_H\ip g{\Hess H(\phi)\,g}
\end{equation}
for all $g\in T$ and $\phi\in\R^N$.
Then for every real $f\in L^2(\pi)$ that is $T^\perp$-invariant, that is,
a function of the orthogonal projection of $\phi$ onto $T$,
\[
\Var_\pi(f)\le\max(m_2,C_H)\sum_{x\in X}\pi\bigl(\Var^\pi_xf\bigr).
\]
In particular, with $\varpi=1/(n^d\max(m_2,C_H))$ and the constant
$C_H$ of~\eqref{eq:rSstable} for the given $T$, the measure $\pi$ satisfies
$\mathrm{BP}(n,\varpi)$ on $L^2(\pi)$ when $T=\R^N$, and on the
$G$-invariant functions $L^2_G$ when $T=G^\perp$. For the partition of unity of
Sec.~\ref{ssec:fffree}, $m_2=S_2^d$, and when $T=\R^N$ every admissible
$C_H$ satisfies $C_H\ge S_2^d$, so that $\varpi=1/(n^dC_H)$.
\end{theorem}

\begin{proof}
Let first $f\in C^\infty(\R^N)$ be $T^\perp$-invariant with bounded
derivatives of all orders. Then $f(\phi+w)=f(\phi)$ for $w\in T^\perp$, so
$\nabla f(\phi)\in T$ for every $\phi$. Let $\Pi$ be the orthogonal
projection of $L^2(\pi;\R^N)$ onto the fields with values in $T$, $(\Pi G)
(\phi)=P_TG(\phi)$. It commutes with $\mathcal L\otimes I$, which acts in
the same way on every component, and with the multiplication by $\Hess H
(\phi)$ by hypothesis; so it commutes with $\mathcal A$ and with $\mathcal
A^{-1}$, and the maximizer $G^*=\mathcal A^{-1}\nabla f=\mathcal A^{-1}\Pi
\nabla f=\Pi G^*$ in~\eqref{eq:rShs} takes values in $T$. Thus
\begin{equation}\label{eq:rSvar}
\Var_\pi(f)=2\pi(\ip{\nabla f}{G^*})-a(G^*,G^*).
\end{equation}

Put $G^{(x)}=\chi_xG^*$, a field supported in $V_x$ (its components outside
$V_x$ vanish) with values in $\R^N$; since $\chi_x$ is a constant diagonal
matrix, $\nabla G^{(x)}_i=\chi_x(i)\nabla G^*_i$, and $G^{(x)}\in\mathcal D$.
As $\sum_x\chi_x=I$, $\sum_xG^{(x)}=G^*$ and
\begin{equation}\label{eq:rSsplit}
2\pi(\ip{\nabla f}{G^*})=\sum_x2\pi\bigl(\ip{\nabla_{V_x}f}{G^{(x)}}\bigr),
\end{equation}
where $\nabla_{V_x}f$ is the restriction of $\nabla f$ to the coordinates of
$V_x$. Fix $x$ and the coordinates outside $V=V_x$. The restriction $g=
G^{(x)}(\cdot,\phi_{V^c})$ is, for $\pi$-almost every exterior, a field on
$\R^V$ in the form domain of $\mathcal A_V$, and the inequality
in~\eqref{eq:rShs} for the conditional law gives
\[
\Var^\pi_xf\ge2\pi_V\bigl(\ip{\nabla_Vf}g\bigr)-a_V(g,g).
\]
Integrating over the exterior,
\begin{equation}\label{eq:rScond}
\pi\bigl(\Var^\pi_xf\bigr)\ge2\pi\bigl(\ip{\nabla_{V_x}f}{G^{(x)}}\bigr)-\pi
\bigl(a_{V_x}(G^{(x)},G^{(x)})\bigr),
\end{equation}
where
\begin{multline*}
\pi\bigl(a_V(G^{(x)},G^{(x)})\bigr)=\sum_{i\in V}\pi\bigl(\lvert\nabla_VG^{(x)}_i
\rvert^2\bigr)\\+\pi\bigl(\ip{G^{(x)}}{\Hess H\,G^{(x)}}\bigr);
\end{multline*}
in the last term we used that $G^{(x)}$ is supported in $V$, so that the
principal submatrix $(\Hess H)_{VV}$ may be replaced by $\Hess H$.

Now $\lvert\nabla_VG^{(x)}_i\rvert^2=\chi_x(i)^2\lvert\nabla_VG^*_i\rvert^2
\le\chi_x(i)^2\lvert\nabla G^*_i\rvert^2$, since the gradient in the
coordinates of $V$ is part of the full gradient. Summing over $x$ and using
$\sum_x\chi_x(i)^2\le m_2$ for the first term and~\eqref{eq:rSstable} at
each $\phi$, with $g=G^*(\phi)\in T$, for the second,
\begin{align*}
\sum_x\pi\bigl(a_{V_x}(G^{(x)},G^{(x)})\bigr)&\le m_2\sum_i\pi\bigl(\lvert\nabla
G^*_i\rvert^2\bigr)\\
&\quad+C_H\,\pi\bigl(\ip{G^*}{\Hess H\,G^*}\bigr)\\
&\le\max(m_2,C_H)\,a(G^*,G^*).
\end{align*}
Put $\lambda=1/\max(m_2,C_H)$, so that $\lambda\sum_x\pi(a_{V_x}(G^{(x)},
G^{(x)}))\le a(G^*,G^*)$. By~\eqref{eq:rSvar} and~\eqref{eq:rSsplit},
\begin{align*}
\Var_\pi(f)&\le\sum_x\Bigl[2\pi\bigl(\ip{\nabla_{V_x}f}{G^{(x)}}\bigr)-
\lambda\,\pi\bigl(a_{V_x}(G^{(x)},G^{(x)})\bigr)\Bigr]\\
&=\lambda^{-1}\sum_x\Bigl[2\pi\bigl(\ip{\nabla_{V_x}f}{\lambda G^{(x)}}
\bigr)\\
&\qquad\qquad-\pi\bigl(a_{V_x}(\lambda G^{(x)},\lambda G^{(x)})\bigr)\Bigr]\\
&\le\lambda^{-1}\sum_x\pi\bigl(\Var^\pi_xf\bigr),
\end{align*}
by~\eqref{eq:rScond} applied to the fields $\lambda G^{(x)}$, which are
again supported in $V_x$ and in the form domains. This is the claim for
smooth $f$ with bounded derivatives.

For general $T^\perp$-invariant $f\in L^2(\pi)$, both sides of the claim are
continuous in $L^2(\pi)$, so it suffices to approximate $f$ in $L^2(\pi)$ by
$T^\perp$-invariant functions of the kind treated. Truncation $f\mapsto\max
(-M,\min(f,M))$ preserves $T^\perp$-invariance and converges in $L^2(\pi)$
as $M\to\infty$. For bounded $f$, convolution with a Gaussian kernel
$\rho_\epsilon$ preserves invariance, since $(f*\rho_\epsilon)(\phi+w)=\int
f(\phi+w-y)\rho_\epsilon(y)\,dy=(f*\rho_\epsilon)(\phi)$ for $w\in T^\perp$,
gives smooth functions with bounded derivatives, and converges to $f$ in
$L^2(\pi)$ as $\epsilon\to0$ because $\pi$ has a density and $f*\rho_\epsilon
\to f$ in $L^2_{\rm loc}(d\phi)$ with $\lvert f*\rho_\epsilon\rvert\le\lVert
f\rVert_\infty$. The statement about $\mathrm{BP}$ is~\eqref{eq:rSBP}. For the partition
of unity of Sec.~\ref{ssec:fffree}, $\sum_x\chi_x(i)^2=S_2^d$ for every $i$
by~\eqref{eq:ffchisums}; and when $T=\R^N$, \eqref{eq:rSstable} is the
operator inequality $\sum_x\chi_x\Hess H(\phi)\chi_x\le C_H\Hess H(\phi)$,
whose trace reads $S_2^d\operatorname{tr}\Hess H(\phi)\le C_H\operatorname{tr}
\Hess H(\phi)$, so $C_H\ge S_2^d$ because $\operatorname{tr}\Hess H(\phi)
\ge Nc>0$. (For $T=G^\perp$ the trace argument is not available; in the
applications below $C_H$ is an explicit bound of the form $S_2^{d-1}(S_2+
\cdots)\ge S_2^d$, and any number at least the best constant
in~\eqref{eq:rSstable} is an admissible $C_H$.)
\end{proof}

\begin{remark}\label{rem:rSwittenremarks}
(a) For a Gaussian, $\Hess H=Q$ is constant and~\eqref{eq:rSstable} is the
stable decomposition hypothesis of Lemma~\ref{lem:fflions} for the specific
decomposition $v_x=\chi_xv$ (with $C_H=C_0(Q;T)$), so the theorem recovers
Theorems~\ref{thm:ffscalar} and~\ref{thm:ffproca} with the same constants.
The chaos decomposition of Theorem~\ref{thm:ffchaos} is replaced by the
Helffer--Sj\"ostrand representation, which plays the same role: it turns a
statement about all functions into a statement about vector fields, on
which the cubes act by multiplication.

(b) When the $V_x$ form a partition of the coordinates into $M$ blocks and
$\chi_x=\mathbf 1_{V_x}$, $m_2=1$ and~\eqref{eq:rSstable} says that the
block-diagonal part $D(\phi)$ of $\Hess H(\phi)$ satisfies $D(\phi)\le C_H
\Hess H(\phi)$; the trace argument gives $C_H\ge1$. The theorem then states
that the random-scan block Gibbs sampler has spectral gap at least
$1/(C_HM)$. Ascolani, Lavenant and Zanella~\cite[Theorem~3.2 and
Corollary~3.7]{ALZ24} prove, for strongly log-concave targets, entropy
contraction at the rate $1/(\kappa^*M)$ and the same Poincar\'e bound with
their coordinate-wise condition number $\kappa^*$, defined through block
smoothness constants $L_m$ with $(\Hess H)_{mm}\le L_mI$ and $\Hess H\ge
(\kappa^*)^{-1}\operatorname{diag}(L_m)$; since $D(\phi)\le\operatorname{diag}
(L_m)$, one has $C_H\le\kappa^*$, and $C_H$ may be taken pointwise in
$\phi$. For Gaussians the bound $1/(C_HM)$ with the best $C_H$ is exact
(Lemma~\ref{lem:fflinear}). The overlapping cubes with the smooth partition
of unity are what makes the constant uniform in $n$: for a partition into
disjoint cubes of side $n$ the constant of the free field is at most
$s^2/(s^2+2dn)$, as the vector $v=\mathbf 1$ in the best constant shows,
because cutting a constant field along hard boundaries costs $2d/n$ per site
against $\mu^2$.

(c) The theorem needs the hypotheses only for the Hessian at every point;
the exterior configuration enters nowhere, so the constant is uniform in
the volume, and in the cube side as soon as $C_H$ is.
\end{remark}

\begin{corollary}[zeroth-order convex perturbations]\label{cor:rSonsite}
Let $\mu_0=\gamma_{Q_{\rm s}}$ be the free scalar field and $\mu=Z^{-1}
e^{-V}\mu_0$ with $V(\phi)=\sum_{y\in\mathbb T}v_y(\phi_y)$, each $v_y\in
C^\infty(\R)$ convex (for instance $v_y(t)=\lambda t^4+\kappa t^2$ with
$\lambda,\kappa\ge0$, or any convex polynomial). Then $\mu$ satisfies
$\mathrm{BP}(n,\varpi)$ on $L^2(\mu)$ with the constant of
Theorem~\ref{thm:ffscalar},
\[
\varpi=\alpha_n^{-d}\frac{s^2}{s^2+10d\rho_n}\ge c_d(s),
\]
uniformly in $L$, in $n$ and in the potentials $v_y$. More generally, if
$R(\phi)\ge0$ is a diagonal matrix for every $\phi$ and $\Hess H=Q_0+R$,
then $C_H\le\max(C_0(Q_0),S_2^d)=C_0(Q_0)$ for the parabolic partition of
unity.
\end{corollary}

\begin{proof}
$H=\frac12\ip\phi{Q_{\rm s}\phi}+V$ has $\Hess H=Q_{\rm s}+R$ with $R=\operatorname{diag}
(v_y''(\phi_y))\ge0$, so $\Hess H\ge\mu^2$. For diagonal $R\ge0$,
$\sum_x\chi_xR\chi_x=R\sum_x\chi_x^2=S_2^dR$, and Lemma~\ref{lem:ffpartition}
gives, for all $g$,
\begin{multline*}
\sum_x\ip{\chi_xg}{(Q_{\rm s}+R)\chi_xg}\le C_0\ip g{Q_{\rm s}g}+S_2^d\ip gRg\\
\le C_0\ip g{(Q_{\rm s}+R)g},
\end{multline*}
with $C_0=S_2^{d-1}(S_2+dD_2/\mu^2)\ge S_2^d$. Theorem~\ref{thm:rSwitten}
with $T=\R^N$ and $C_H=C_0$ gives $\varpi\ge1/(n^dC_0)$, which is the
constant of Theorem~\ref{thm:ffscalar}. The integrability of $e^{-H}$ and
the moments needed for~\eqref{eq:rShs} hold because a convex $V$ is bounded
below by an affine function.
\end{proof}

\begin{corollary}[two-sided Hessian bounds]\label{cor:rStwosided}
Let $Q_0$ be positive definite, $K\ge1$, and let $\pi=Z^{-1}e^{-H}d\phi$ with
$H\in C^\infty$ and
\[
Q_0\le\Hess H(\phi)\le KQ_0\qquad(\phi\in\R^N).
\]
Then $\pi$ satisfies $\mathrm{BP}(n,\varpi)$ on $L^2(\pi)$ with $\varpi=1/
(n^dKC_0(Q_0))$. If moreover $\Hess H(\phi)$ commutes with the projection
onto a subspace $G$ for every $\phi$, then $\mathrm{BP}(n,\varpi_G)$ holds on
$L^2_G$ with
\[
\varpi_G=\frac1{n^d\max\bigl(S_2^d,\,KC_0(Q_0;G^\perp)\bigr)}.
\]
For the free scalar and Proca
fields, with the explicit bounds of Lemmas~\ref{lem:ffpartition}
and~\ref{lem:ffprocadecomp} in place of $C_0$, this gives the constants of
Theorems~\ref{thm:ffscalar} and~\ref{thm:ffproca} divided by $K$.
\end{corollary}

\begin{proof}
For $g\in T$ ($T=\R^N$ or $G^\perp$), the upper bound at the vectors
$\chi_xg$, which need not lie in $T$, and the lower bound at $g$ give
\begin{multline*}
\sum_x\ip{\chi_xg}{\Hess H\chi_xg}\le K\sum_x\ip{\chi_xg}{Q_0\chi_xg}\\
\le KC_0(Q_0;T)\ip g{Q_0g}\le KC_0(Q_0;T)\ip g{\Hess Hg},
\end{multline*}
and $\Hess H\ge Q_0\ge\lambda_{\min}(Q_0)>0$. Apply
Theorem~\ref{thm:rSwitten} with $C_H=KC_0(Q_0;T)$, or with $C_H=K\bar C$
for any upper bound $\bar C$ of $C_0(Q_0;T)$; for $T=\R^N$, $C_0(Q_0)\ge
S_2^d$ by the trace argument in that proof, and the explicit bounds of the
two lemmas are of the form $S_2^{d-1}(S_2+\cdots)\ge S_2^d$.
\end{proof}

\begin{corollary}[convex plaquette perturbations of the Proca field]
\label{cor:rSplaquette}
Let $d\ge2$, $\mu_0=\gamma_{Q_{\rm P}}$ the free Proca field, and $\mu=Z^{-1}
e^{-V}\mu_0$ with $V(A)=\sum_{p\in P}v_p\bigl((d_1A)(p)\bigr)$, each $v_p\in
C^\infty(\R)$ convex with $v_p''\le K$. Then $V$ is gauge-invariant and $\mu$
satisfies $\mathrm{BP}(n,\varpi)$ on the gauge-invariant functions with
\[
\varpi=\frac1{1+K}\,\alpha_n^{-d}\frac{s^2}{s^2+10d\rho_n}\ge\frac{c_d(s)}{1+K},
\]
uniformly in $L$ and $n$.
\end{corollary}

\begin{proof}
$\Hess V(A)=d_1^{\mathsf T}D(A)d_1$ with $D=\operatorname{diag}(v_p''((d_1A)(p)))$,
$0\le D\le K$; it vanishes on $G=d_0\R^{\mathbb T}$ and maps into $G^\perp$,
so it commutes with the projection onto $G$, as does $Q_{\rm P}$ by
Lemma~\ref{lem:ffhodge}; and $\Hess H=Q_{\rm P}+\Hess V\ge\mu^2$. For
transverse $g$,
\begin{multline*}
\sum_x\ip{\chi_xg}{\Hess V\chi_xg}=\sum_x\lVert D^{1/2}d_1(\chi_xg)\rVert^2\\
\le K\sum_x\lVert d_1(\chi_xg)\rVert^2\le K\sum_x\ip{\chi_xg}{Q_{\rm P}\chi_xg},
\end{multline*}
so by Lemma~\ref{lem:ffprocadecomp} (second bound)
\begin{multline*}
\sum_x\ip{\chi_xg}{\Hess H\chi_xg}\le(1+K)C_0\ip g{Q_{\rm P}g}\\
\le(1+K)C_0\ip g{\Hess Hg},
\end{multline*}
$C_0=S_2^{d-1}(S_2+dD_2/\mu^2)\ge S_2^d=m_2$. Theorem~\ref{thm:rSwitten}
with $T=G^\perp$ and $C_H=(1+K)C_0$ gives the claim, since $n^dC_0=
\alpha_n^d(1+10d\rho_ns^{-2})$.
\end{proof}

The plaquette weight of the Wilson action, $v(\theta)=\beta(1-\cos\theta)$,
has $v''=\beta\cos\theta$ and is excluded by the corollary twice: it is not
convex, and the Proca mass is absent. Both failures are structural, and
Sec.~\ref{sssec:rSgauge} makes them precise.

\subsubsection{Bounded perturbations}\label{sssec:rSbounded}

\begin{proposition}[global Holley--Stroock bound]\label{prop:rSglobalHS}
Let $\mu_0$ satisfy $\mathrm{BP}(n,\varpi_0)$ on $L^2(\mu_0)$ and let $\mu=
Z^{-1}e^{-W}\mu_0$ with $W$ bounded. Let $\omega=\sup W-\inf W$ and
\[
\omega_{\rm loc}=\max_x\sup_{\phi_{V_x^c}}\operatorname{osc}_{\phi_{V_x}}W
\le\omega,
\]
the largest oscillation of $W$ over a cube at fixed exterior. Then $\mu$
satisfies $\mathrm{BP}(n,e^{-2\omega-\omega_{\rm loc}}\varpi_0)$ on
$L^2(\mu)$.
\end{proposition}

\begin{proof}
The density $d\mu/d\mu_0=e^{-W}/\mu_0(e^{-W})$ lies in $[e^{-\omega},e^{
\omega}]$, and the conditional densities given the exterior of $V_x$ are
$e^{-W}/\mu_{0,x}(e^{-W})$, in $[e^{-\omega_{\rm loc}},e^{\omega_{\rm
loc}}]$. Since $\Var_\nu(f)=\inf_c\nu((f-c)^2)$, a density in $[e^{-\omega},
e^\omega]$ changes variances by at most the factor $e^{\pm\omega}$. Hence
\begin{align*}
\Var_\mu f&\le e^\omega\Var_{\mu_0}f\le\frac{e^\omega}{\varpi_0n^d}\sum_x\mu_0
\bigl(\Var^{\mu_0}_xf\bigr)\\
&\le\frac{e^{2\omega}}{\varpi_0n^d}\sum_x\mu\bigl(\Var^{\mu_0}_xf\bigr)
\le\frac{e^{2\omega+\omega_{\rm loc}}}{\varpi_0n^d}\sum_x\mu\bigl(\Var^\mu_xf\bigr).
\end{align*}
\end{proof}

The global oscillation $\omega$ grows with the volume for a local
perturbation $W=\sum_pW_p$, and the proposition is useless on large tori.
Only the local oscillation $\omega_{\rm loc}$, of order $n^d\max_p
\operatorname{osc}W_p$, is uniform in the volume, and one may ask for a bound
$\varpi(\mu)\ge F(d,n,\omega_{\rm loc})\varpi(\mu_0)$. No such bound exists.

\begin{proposition}[product measures]\label{prop:rSproduct}
Every product probability measure on a finite product of measurable spaces
indexed by the coordinates satisfies $\mathrm{BP}(n,1)$ on $L^2$, for every
family $(V_x)$ in which every coordinate lies in at least $n^d$ sets.
\end{proposition}

\begin{proof}
Let $\mathcal E_i$ be the expectation in the $i$-th coordinate. For $f\in
L^2$ and $S\subset\{1,\dots,N\}$ put $f_S=\prod_{i\in S}(I-\mathcal E_i)
\prod_{i\notin S}\mathcal E_if$; the operators commute, $f=\sum_Sf_S$, the
$f_S$ are mutually orthogonal, $f_\emptyset=\E f$ and $\mathcal P_xf_S=f_S$
if $S\cap V_x=\emptyset$ while $\mathcal P_xf_S=0$ otherwise (it contains a
factor $\mathcal E_i(I-\mathcal E_i)=0$). So $\lVert f-\mathcal P_xf\rVert^2
=\sum_{S\cap V_x\ne\emptyset}\lVert f_S\rVert^2$ and $\sum_x\lVert f-\mathcal
P_xf\rVert^2=\sum_{S\ne\emptyset}\lVert f_S\rVert^2\,\#\{x:V_x\cap S\ne
\emptyset\}\ge n^d\Var f$.
\end{proof}

\begin{proposition}[no local Holley--Stroock bound]\label{prop:rSising}
Let $d\ge3$, $L\ge n+1$ even, and let $\mu_\beta$ be the Ising measure on
$\{\pm1\}^{\mathbb T}$, $\mu_\beta(\sigma)\propto\exp(\beta\sum_{\langle yz
\rangle}\sigma_y\sigma_z)$, the sum over the nearest-neighbour pairs of the
torus. Then $\mu_\beta=Z^{-1}e^{-W}\mu_0$ with $\mu_0$ the uniform measure,
which satisfies $\mathrm{BP}(n,1)$, and $W=-\beta\sum_{\langle yz\rangle}
\sigma_y\sigma_z$, whose local oscillation is at most $2\beta d(n^d+n^{d-1})$.
There is $\beta_1(d)<\infty$ such that for $\beta\ge\beta_1(d)$ and every
even $L$,
\[
\varpi(\mu_\beta)\le\frac{2n^d}{L^d};
\]
one may take $\beta_1(3)=6.7$ and $\beta_1(4)=8.6$. Consequently there is no
function $F$ with $\varpi(e^{-W}\mu_0)\ge F(d,n,\omega_{\rm loc}(W))\,\varpi
(\mu_0)$ on all tori: the cube condition is not stable under bounded local
perturbations with constants uniform in the volume.
\end{proposition}

\begin{proof}
The uniform measure is a product measure, so Proposition~\ref{prop:rSproduct}
gives $\mathrm{BP}(n,1)$; a cube meets at most $d(n^d+n^{d-1})$ bonds, each
contributing oscillation $2\beta$. Let $M=\sum_y\sigma_y$ and take $f=M$
in~\eqref{eq:rSBP}. Given the exterior of $V_x$, $M$ varies only through
$\sum_{y\in V_x}\sigma_y\in[-n^d,n^d]$, so $\Var_x^{\mu_\beta}M\le n^{2d}$
and $n^{-d}\sum_x\mu_\beta(\Var_xM)\le L^dn^d$. For the variance we use the
infrared bound of Fr\"ohlich, Simon and Spencer~\cite{FSS76}, in the torus
form of~\cite[Theorem~3.3 and Corollary~3.4]{Bis09}: for $L$ even,
$\hat\sigma(k)=\sum_y\sigma_ye^{ik\cdot y}$ and $k\in\mathbb T^*\setminus\{0\}$,
$\mathbb T^*=\frac{2\pi}L\{0,\dots,L-1\}^d$,
\begin{equation}\label{eq:rSirb}
\mu_\beta\bigl(\lvert\hat\sigma(k)\rvert^2\bigr)\le\frac{L^d}{2\beta D(k)},
\qquad D(k)=\sum_{j=1}^d(1-\cos k_j).
\end{equation}
The constant depends on the normalisation of the Hamiltonian: \cite{Bis09}
writes the bound for $-\frac1{2d}\sum_{\langle yz\rangle}\sigma_y\sigma_z$ at
inverse temperature $2d\beta$, with $1-\hat J(k)=D(k)/d$, and in the present
normalisation~\eqref{eq:rSirb} is the second-order expansion at $h=\cos(k
\cdot y)$ and $h=\sin(k\cdot y)$ of Gaussian domination, $\mu_\beta\bigl(
e^{\beta\ip\sigma{-\Delta h}}\bigr)\le e^{\frac\beta2\ip h{-\Delta h}}$, for
which $\mu_\beta(\ip\sigma h^2)\le\lVert h\rVert^2/(2\beta D(k))$ when $h$
is an eigenvector of $-\Delta$ with eigenvalue $2D(k)$. By Parseval,
$\sum_{k\in\mathbb T^*}\lvert\hat\sigma(k)\rvert^2=L^d\sum_y\sigma_y^2=L^{2d}$,
and $\hat\sigma(0)=M$, so
\begin{gather*}
\mu_\beta(M^2)\ge L^{2d}\Bigl(1-\frac{G_L}{2\beta}\Bigr),\qquad
G_L=\frac1{L^d}\sum_{k\in\mathbb T^*\setminus\{0\}}\frac1{D(k)}.
\end{gather*}
We bound $G_L$ uniformly in $L$. For $k\in\mathbb T^*$ identify $k$ with
its representative in $(-\pi,\pi]^d$; then $1-\cos k_j\ge2k_j^2/\pi^2$, so
$D(k)\ge2\lvert k\rvert^2/\pi^2$. For $k\ne0$ let $B_k$ be the cube of side
$2\pi/L$ centred at $k$; the $B_k$ are disjoint, contained in $[-2\pi,2\pi]^d$,
and for $k'\in B_k$, $\lvert k'\rvert\le\lvert k\rvert+\sqrt d\,\pi/L\le(1+
\sqrt d/2)\lvert k\rvert$ because $\lvert k\rvert\ge2\pi/L$. Hence
\[
G_L\le\frac{\pi^2}2\Bigl(1+\frac{\sqrt d}2\Bigr)^2\frac1{(2\pi)^d}\int_{[-2
\pi,2\pi]^d}\frac{dk'}{\lvert k'\rvert^2}=:I_d,
\]
which is finite for $d\ge3$; with $k'=2\pi u$, $I_d=\frac18(1+\frac{\sqrt d}
2)^2\int_{[-1,1]^d}\lvert u\rvert^{-2}du$, and the integral is $15.35$ for
$d=3$ and $17.15$ for $d=4$, so $I_3<6.7$ and $I_4<8.6$. With $\beta_1(d)=
I_d$, for $\beta\ge\beta_1$ we get $G_L/(2\beta)\le\frac12$ and
$\Var_{\mu_\beta}(M)=\mu_\beta(M^2)\ge L^{2d}/2$, using $\mu_\beta(M)=0$ by
the spin-flip symmetry. Then~\eqref{eq:rSBP} with $f=M$ forces $\varpi\le
L^dn^d/(L^{2d}/2)=2n^d/L^d$. The last claim follows by letting $L\to\infty$
at fixed $n$ and $\beta$, since $\omega_{\rm loc}(W)$ and $\varpi(\mu_0)=1$
do not depend on $L$.
\end{proof}

What survives for bounded perturbations is smallness measured against the
Hessian, not against the oscillation: if $W\in C^\infty$ satisfies $-(1-
\epsilon)Q_0\le\Hess W\le KQ_0$, then $\epsilon Q_0\le\Hess(\tfrac12\ip\phi
{Q_0\phi}+W)\le(1+K)Q_0$ and Corollary~\ref{cor:rStwosided} gives
$\mathrm{BP}(n,\epsilon/((1+K)n^dC_0(Q_0)))$. For the massive field $Q_0\ge
\mu^2=(ma)^2$, so a bounded perturbation with $\lVert\Hess W\rVert_{\rm op}
\le\delta$ is admissible by this route only when $\delta<(ma)^2$: in the
continuum tuning the perturbation must be of order $a^2$ in lattice units,
which is the lattice form of the statement that such a perturbation is
irrelevant. Perturbations of order one in lattice units with small
oscillation per term are the domain of the Dobrushin--Shlosman
condition~\cite{Dob68,DS85}, which the paper uses in
Proposition~\ref{prop:avgsup}; Proposition~\ref{prop:rSising} shows that
smallness of the oscillation per cube cannot replace it.

\subsubsection{The lattice gauge field}\label{sssec:rSgauge}

We now place the Wilson measure against the hypotheses of
Theorem~\ref{thm:rSwitten}. For definiteness we take the group $\U(1)$, with
link variables $\theta\in(\R/2\pi\Z)^E$, plaquette angles $\theta_p=(d_1
\theta)(p)$ and action
\[
H_\beta(\theta)=\beta\sum_{p\in P}(1-\cos\theta_p),\qquad\mu_\beta\propto
e^{-H_\beta}\prod_ed\theta_e ,
\]
on the torus of Sec.~\ref{ssec:fffree} with $d\ge2$; the remarks on $\SU(N)$
follow the statements. A \emph{linear gauge slice} is a subspace $S\subset
\R^E$ with $S\oplus G=\R^E$, $G=d_0\R^{\mathbb T}$ the pure gauge directions:
the Landau slice $S=G^\perp=\ker\delta_0$ and the axial slices, on which the
links of a maximal tree are set to zero (the tree links are in bijection
with $G$, since gauge transformations modulo constants are determined by
their values on a spanning tree), are examples. Let $F=\ker d_1\cap\ker
\delta_0$ be the flat transverse fields.

The two consequences of the Helffer--Sj\"ostrand representation isolated
after~\eqref{eq:rShs} survive the loss of convexity on the torus. The
generator $\mathcal L=-\Delta+\ip{\nabla H_\beta}\nabla$ on $L^2(\mu_\beta)$
has compact resolvent and kernel the constants, since the configuration
space is compact and the density is smooth and positive, so $u_f=\mathcal
L^{-1}(f-\mu_\beta(f))$ exists for every smooth $f$; the form
\[
a(G,G)=\sum_{e\in E}\mu_\beta\bigl(\lvert\nabla G_e\rvert^2\bigr)+\mu_\beta
\bigl(\ip G{\Hess H_\beta\,G}\bigr)
\]
is the quadratic form of the Witten Laplacian on one-forms of the flat torus
with the weight $e^{-H_\beta}$, which is $d_H^*d_H+d_Hd_H^*$ by the Bochner
formula and hence nonnegative for every $H_\beta$; and the intertwining
relation gives $a(\nabla u_f,G)=\mu_\beta(\ip{\nabla f}G)$ for every $G$ in
the form domain. Hence, with $G_f:=\nabla u_f$,
\begin{align}\label{eq:rShstorus}
\Var_{\mu_\beta}f&=2\mu_\beta(\ip{\nabla f}{G_f})-a(G_f,G_f)\notag\\
&\ge2\mu_\beta(\ip{\nabla f}G)-a(G,G)
\end{align}
for every such $G$, and the same holds for the conditional laws on the cubes,
whose Hamiltonians are the restrictions of $H_\beta$, with the forms $a_{V_x}$
defined as before. What is lost is only the lower bound $a\ge c$ and the
invertibility of $\mathcal A$: by Proposition~\ref{prop:rSwilson}(iv) below
the constant flat fields lie in its kernel.

\begin{proposition}[the Hessian of the Wilson action]\label{prop:rSwilson}
Let $d\ge2$, $L\ge3$ and $\beta>0$.
\begin{enumerate}[label=\textup{(\roman*)},leftmargin=*]
\item The Hessian is $\Hess H_\beta(\theta)=\beta\,d_1^{\mathsf T}D_\theta
d_1$ with $D_\theta=\operatorname{diag}(\cos\theta_p)$; it vanishes on $G$
and on $F$ at every $\theta$, and $F$ is the $d$-dimensional space of the
constant fields $A(y,i)=c_i$.
\item $\Hess H_\beta(0)=\beta\delta_1d_1$, and $\ker\Hess H_\beta(0)=G\oplus
F$. For every linear gauge slice $S$ the restriction of $\Hess H_\beta(0)$ to
$S$ has a kernel of dimension at least $d$. Hence there is no positive
definite $Q_0$ on $S$, nor on $G^\perp$, with $\Hess H_\beta\ge Q_0$ on a
neighbourhood of $0$, for any $\beta$.
\item For every $\beta>0$ the set of $\theta$ at which $\Hess H_\beta(\theta)$
has a negative eigenvalue on $G^\perp$ is open and has positive $\mu_\beta$
measure; it contains every $\theta$ at which the $2(d-1)$ plaquettes through
some link $e$ have $\sum_{p\ni e}\cos\theta_p<0$.
\item For every $\beta>0$, $c_\beta=\mu_\beta(\cos\theta_p)>0$, and for the
constant field $h$ in direction $1$, $h(y,i)=\delta_{i1}$, the Witten form of
$\mu_\beta$ and its cube restrictions satisfy
\begin{gather*}
a(h,h)=\mu_\beta\bigl(\ip h{\Hess H_\beta h}\bigr)=0,\\
\sum_x\mu_\beta\bigl(a_{V_x}(\chi_xh,\chi_xh)\bigr)=\beta c_\beta(d-1)S_2^{d-1}D_2L^d>0 .
\end{gather*}
\end{enumerate}
\end{proposition}

\begin{proof}
(i) $H_\beta$ is a function of $d_1\theta$, so $\nabla H_\beta=\beta d_1^{
\mathsf T}\sin(d_1\theta)$ and $\Hess H_\beta=\beta d_1^{\mathsf T}\operatorname
{diag}(\cos\theta_p)d_1$, which vanishes on $\ker d_1\supset G\cup F$. By the
Hodge identity of Lemma~\ref{lem:ffhodge}, $A\in F$ has $\sum_i\lVert\nabla
A_i\rVert^2=0$, so each $A_i$ is constant on the connected torus; conversely
constant fields have $d_1A=0=\delta_0A$. So $F$ is the space of constants,
of dimension $d$.

(ii) At $\theta=0$ all cosines are $1$. $\ker\delta_1d_1=\ker d_1$ and $\ker
d_1=G\oplus F$: indeed $G\subset\ker d_1$, $F\subset\ker d_1$, $G\perp F$,
and for $A\in\ker d_1$ the Hodge decomposition $A=P_GA+P_{G^\perp}A$ has
$d_1P_{G^\perp}A=d_1A-d_1P_GA=0$ and $\delta_0P_{G^\perp}A=0$, so $P_{G^\perp}A
\in F$. For a slice $S$, $\dim(S\cap\ker d_1)=\dim S+\dim\ker d_1-\dim(S+
\ker d_1)\ge(N-\dim G)+(\dim G+d)-N=d$, as $S+\ker d_1\supset S+G=\R^E$; and
$S\cap\ker d_1\subset\ker\Hess H_\beta(0)|_S$. A bound $\Hess H_\beta\ge Q_0$
near $0$ would give $Q_0\le0$ on this kernel; on $G^\perp$ the kernel contains
$F\ne\{0\}$.

(iii) For a link $e$ and $\theta$ as stated, $\ip{\delta_e}{\Hess H_\beta
\delta_e}=\beta\sum_{p\ni e}\cos\theta_p<0$, each plaquette through $e$
containing $e$ with coefficient $\pm1$. Since $\Hess H_\beta$ is symmetric and
vanishes on $G$, $\ip{P_{G^\perp}\delta_e}{\Hess H_\beta P_{G^\perp}\delta_e}=
\ip{\delta_e}{\Hess H_\beta\delta_e}<0$, and $P_{G^\perp}\delta_e\ne0$ because
$\delta_e\notin G$ ($d_1\delta_e\ne0$). The set is open as $\Hess H_\beta$ is
continuous in $\theta$, nonempty (take $\theta_e=\pi$ and all other links
zero, so that every plaquette through $e$ has angle $\pm\pi$), and the
density of $\mu_\beta$ is positive everywhere.

(iv) $\log Z(\beta)=\log\int e^{\beta\sum_p\cos\theta_p}d\theta$ is convex
in $\beta$ with derivative $\lvert P\rvert\,\mu_\beta(\cos\theta_p)$ (all
plaquettes are equivalent under the translations and the hypercubic group,
and $\mu_\beta$ is invariant), and its derivative at $0$ is $\lvert P\rvert
\int\cos\theta_p\,d\theta=0$; the second derivative is $\Var_{\mu_\beta}(\sum_p
\cos\theta_p)>0$ since the plaquette angles are not constant, so the
derivative is strictly increasing and $c_\beta>0$ for $\beta>0$. The constant
vector field $h$ has $\nabla h=0$, so $a(h,h)=\mu_\beta(\ip h{\Hess H_\beta
h})=\beta\mu_\beta(\sum_p\cos\theta_p(d_1h)_p^2)=0$ by (i), and likewise
$a_{V_x}(\chi_xh,\chi_xh)=\beta\sum_p\cos\theta_p\bigl(d_1(\chi_xh)\bigr)_p^2$,
whose expectation is $\beta c_\beta\lVert d_1(\chi_xh)\rVert^2$. Now
$(\chi_xh)(y,i)=\chi_x(y)\delta_{i1}$, so $d_1(\chi_xh)(y,1,j)=\chi_x(y)-
\chi_x(y+e_j)=-\nabla_j\chi_x(y)$ for $j\ne1$ and $d_1(\chi_xh)(y,i,j)=0$ for
$i,j\ne1$. By~\eqref{eq:ffchisums}, $\sum_x(\nabla_j\chi_x(y))^2=2S_2^d-
2S_2^{d-1}S_{11}=S_2^{d-1}D_2$, and summing over the $(d-1)L^d$ pairs $(y,j)$
gives the claim.
\end{proof}

For $\SU(N)$ the link variables are $U_e=e^{X_e}$, $X_e\in\mathfrak{su}(N)$,
and the Hessian of $\beta\sum_p(1-\frac1N\ReTr U_p)$ at the trivial
configuration is $\frac\beta{N}\delta_1d_1$ on each of the $N^2-1$
components with respect to a basis of $\mathfrak{su}(N)$ orthonormal for
$-\Tr(XY)$, so (ii) holds verbatim with $\dim F=d(N^2-1)$; the statement at
the trivial configuration is a statement about the Maxwell form and does not
see the group. The paper's Sec.~\ref{sec:curvature} records the analogue of
(iii) for $\SU(2)$: the Bakry--\'Emery tensor is negative at large fields
once $\beta_W>\frac13$.

\begin{corollary}[the positive theorems do not apply to the Wilson measure]
\label{cor:rSnoapply}
For every $\beta>0$, every $d\ge2$ and every linear gauge slice $S$, the
Wilson action restricted to $S$ violates each hypothesis of
Theorem~\ref{thm:rSwitten} and of Corollaries~\ref{cor:rStwosided}
and~\ref{cor:rSplaquette}: it is not convex on an open set of positive
measure, and even where it is convex its Hessian has no positive lower
bound, since $S\cap\ker d_1$, of dimension at least $d$, lies in the kernel
of $\Hess H_\beta(\theta)|_S$ at every $\theta$. Moreover, for the
configuration-independent test field $h$ of
Proposition~\ref{prop:rSwilson}(iv) the inequality
\begin{equation}\label{eq:rSstar}
a(G,G)\ge\lambda\sum_x\mu_\beta\bigl(a_{V_x}(\chi_xG,\chi_xG)\bigr)
\end{equation}
fails for every $\lambda>0$, so that~\eqref{eq:rSstar} cannot hold for all
fields $G$, for any partition of unity and any $\beta$.
\end{corollary}

\begin{proof}
Proposition~\ref{prop:rSwilson}(ii), (iii) and (iv); for the kernel at every
$\theta$ use that $\ker d_1\subset\ker\Hess H_\beta(\theta)$ by (i) and the
dimension count of (ii). For the last claim note that~\eqref{eq:rSstar} at
$G=h$ reads $0\ge\lambda\beta c_\beta(d-1)S_2^{d-1}D_2L^d>0$; any other
partition of unity with $\sum_x\chi_x=I$ and $\chi_x$ supported in $V_x$
cuts the constant field $h$ along the boundary of every cube, so that
$\sum_x\lVert d_1(\chi_xh)\rVert^2>0$ as well.
\end{proof}

Inequality~\eqref{eq:rSstar} for all $G$ is sufficient for the cube
condition, by the proof of Theorem~\ref{thm:rSwitten} with
\eqref{eq:rShstorus} in place of~\eqref{eq:rShs}, and it is not necessary:
the field $h$ is the gradient of no function on the torus, since $\nabla u=h$
would force $u(\theta+th)=u(\theta)+t\lvert h\rvert^2$, and the generator part
of the Witten form, $\sum_e\mu_\beta(\lvert\nabla G_e\rvert^2)$, vanishes on
configuration-independent fields. So the corollary does not refute the cube
condition for the Wilson measure. It locates the estimate that any proof
along these lines must supply.

\begin{proposition}[the quadratic approximation has no mass]
\label{prop:rSquadratic}
The quadratic approximation of the Wilson action at the trivial
configuration is the quadratic form $\beta\delta_1d_1$, the free Maxwell
action; it is degenerate on $F$, so it defines no Gaussian measure, even on a
gauge slice. For its massive regularization, the Gaussian with precision
$\beta(\delta_1d_1+\mu^2)$, Proposition~\ref{prop:ffupper}(ii) gives, on
gauge-invariant functions and for every $\beta$,
\[
\varpi\le\frac{(n\mu)^2}{1+(n\mu)^2},
\]
which tends to $0$ as $\mu\to0$. Consequently, along a continuum tuning with
$n_ka_k=\ell$, a Gaussian approximation of this form satisfies $\mathrm{BP}
(n_k,\varpi)$ uniformly in $k$ only if its mass $\mu_k$ satisfies $\mu_k\ge
\sqrt{\varpi/(1-\varpi)}\,a_k/\ell$, that is, only if it carries a physical
mass $m\ge\sqrt{\varpi/(1-\varpi)}/\ell$; and the Wilson action has no bare
mass term.
\end{proposition}

\begin{proof}
The precision $\beta(\delta_1d_1+\mu^2)$ is a multiple of $Q_{\rm P}$, and
$\varpi$ is homogeneous of degree zero in the precision (proof of
Theorem~\ref{thm:rSnonmono}), so Proposition~\ref{prop:ffupper}(ii) applies.
Solving $(n\mu)^2/(1+(n\mu)^2)\ge\varpi$ gives $n\mu\ge\sqrt{\varpi/(1-
\varpi)}$, and $n_k\mu_k=\ell\mu_k/a_k$.
\end{proof}

\begin{remark}[where a dynamical mass would have to enter]\label{rem:rSmissing}
Theorem~\ref{thm:avgtransfer} shows that $\mathrm{BP}(n_k,\varpi)$ along the
tuning gives a mass gap $m_*=16\varpi/(625e^2\ell)$, and
Proposition~\ref{prop:rSquadratic} with Theorem~\ref{thm:ffproca} shows that
in the Gaussian class with precision $\beta(\delta_1d_1+\mu^2)$ the constant
$\varpi$ is of the order of the squared mass in units of the cube when that
mass is at most one. For the Wilson measure the constant therefore has to be
produced by what is absent from the quadratic approximation.
Proposition~\ref{prop:rSwilson} says where it cannot come from: not from a
Hessian lower bound (there is none, in any gauge, at any $\beta$), and not
from the potential term of the Witten form on flat directions, which vanishes
identically. In the Helffer--Sj\"ostrand representation~\eqref{eq:rShstorus}
the cube condition on gauge-invariant functions is implied by the following
statement, which is the proof of Theorem~\ref{thm:rSwitten} with
configuration-dependent decompositions: there is $\lambda>0$ such that for
every smooth gauge-invariant $f$ the field $G_f=\nabla\mathcal L^{-1}(f-
\mu_\beta(f))$ admits a decomposition $G_f=\sum_xG_f^{(x)}$ into fields in
the form domain with $G_f^{(x)}(\theta)$ supported in $V_x$ for every
$\theta$, such that
\begin{equation}\label{eq:rSmissing}
\sum_x\mu_\beta\bigl(a_{V_x}(G_f^{(x)},G_f^{(x)})\bigr)\le\lambda^{-1}a(G_f,
G_f);
\end{equation}
then $\mathrm{BP}(n,\lambda/n^d)$ holds, since that proof used
\eqref{eq:rSstable} and $m_2$ only to produce such a decomposition.
Corollary~\ref{cor:rSnoapply} shows that the decomposition $G^{(x)}=\chi_xG$,
which under the hypotheses of Theorem~\ref{thm:rSwitten} works for every
field at once, cannot be controlled for the Wilson measure by a bound of the
form~\eqref{eq:rSstar} valid for all fields: a field that is constant in
$\theta$ and flat in space has $a(G,G)=0$ and positive cube energy. Whether
\eqref{eq:rSmissing} holds with $G^{(x)}_f=\chi_xG_f$ for the fields $G_f$
that actually occur is not decided here; these fields are orthogonal in
$L^2(\mu_\beta;\R^E)$ to the constant flat fields, since $\mu_\beta(\ip{\nabla
u}h)=\mu_\beta(u\,\partial_hH_\beta)=0$ for flat $h$, but not pointwise in
$\theta$. If the flat part of $G_f$ is what the cut costs, the term $\sum_e
\mu_\beta(\lvert\nabla_\theta G_e\rvert^2)$ of the Witten form, the
fluctuation of the field as a function of the configuration, has to pay for
the curvature created at the cube boundaries. By
Proposition~\ref{prop:rSwilson}(iv) that curvature costs $\beta c_\beta(d-1)
S_2^{d-1}D_2$ per unit of $\lvert h\rvert^2$ for the flat field $h$. In the
Gaussian class the same cut is paid by the mass term, which contributes
$\beta\mu^2S_2^d$ per unit of $\lvert h\rvert^2$ to the left side
of~\eqref{eq:rSstable} for the precision $\beta(\delta_1d_1+\mu^2)$, and the
ratio of the two rates is $(d-1)c_\beta D_2/(\mu^2S_2)=(d-1)c_\beta\beta_n/
(\alpha_ns^2)$, with $\beta_n/\alpha_n\to10$ and $s=n\mu=m\ell$ along the
tuning; so the mass term dominates the cut when $s^2\ge10\rho_n(d-1)c_\beta$,
that is $m\ell\gtrsim\sqrt{10(d-1)c_\beta}$, a requirement of the same form
as that of Proposition~\ref{prop:rSquadratic}. This is a form, in the present
framework, of the heuristic statement that the mass must be dynamical: what
is missing is a lower bound on the Dirichlet form of the Langevin generator of
$\mu_\beta$ on the components of the one-forms $G_f$, uniform in the volume
and in the coupling along the tuning, which on the flat directions plays the
role of the mass term $(ma)^2$ of the Proca field. No such bound is known;
the published control of the lattice measure is either at strong coupling,
where the Bakry--\'Emery tensor is positive and the mass is of order $1/a$
(Sec.~\ref{sec:curvature}, \cite{SZZ23}), or in the small-field region after
a large-field split~\cite{Bal89b}, where the quadratic approximation is
massless. Both are consistent with Proposition~\ref{prop:rSquadratic}: a
lattice-unit mass gives $(n\mu)^2\to\infty$, a massless approximation gives
$0$.
\end{remark}
